\chapter*{Abstract}
\addcontentsline{toc}{chapter}{Abstract}

Kotlin/Wasm is a compilation target for the Kotlin programming language that compiles Kotlin code to
WebAssembly. It is used, for example, by the Compose Multiplatform library for its Web target.
Kotlin Coroutines are central to Kotlin/Wasm applications, yet the current implementation is based on
state-machine transformation: each suspend function is compiled into a state machine that saves and
restores local variables at each suspension point. This approach adds code that increases binary size
and introduces runtime overhead.

The WebAssembly Stack-switching proposal provides native primitives for suspending and resuming
execution contexts. This thesis presents an alternative implementation of Kotlin Coroutines for
Kotlin/Wasm based on these primitives, replacing the state-machine transformation with direct
stack-switching instructions.

A key technical challenge addressed in this work is the mismatch between Kotlin's coroutine model,
where a continuation object must be available before a function suspends, and the Wasm Stack-switching
model, where the continuation only becomes available after suspension. The thesis describes two
implementations that resolve this mismatch.

The evaluation shows that the Stack-switching implementation reduces binary size by approximately
10.8\% in production builds and 2.3\% in development builds for a small Compose Multiplatform
project. However, performance benchmarks on V8 show a noticeable regression: the most realistic
benchmark cases are around 2 times slower than the baseline, and cases with intensive suspension
chaining are 4--4.5 times slower without an additional V8 flag, or around 1.5 times slower with it.
The performance regression is likely related to the current state of the V8 Stack-switching
implementation, which is still under active development.

